\chapter{Судалгааны үр дүн}

\section{Туршилтын үр дүн ба үнэлгээ}

Энэ бүлэгт боловсруулсан ирц бүртгэлийн системийн туршилтын үр дүн, үнэлгээний протокол, хэмжүүрүүдийн хэрэглээ, алдааны шинжилгээ болон практик хэрэглээний хязгаарлалтыг нэгтгэн тайлагнана.

\subsection{Туршилтын өгөгдлийн тайлбар}
Системийн үнэлгээ нь анги доторх бодит орчны зураг болон ирцийн CSV бүртгэлд тулгуурлан хийгдэнэ. Өгөгдлийн эх үүсвэр нь автоматаар үүссэн сешн файл, баталгаажсан ирцийн файл, мөн аннотацтай зурагнуудаас бүрдэнэ. Үнэлгээ хийхдээ дараах үндсэн нэгжүүдийг ашиглана:

\begin{itemize}
	\item Мөрийн түвшин: \texttt{attendance/sessions/*.csv} файл дахь нэг мөр нь \texttt{хугацааны тэмдэглэгээ(timestamp), файлын нэр(filename), таньсан сурагчийн индекс(face\_index), сурагчийн оюутны код болон нэр(label), accuracy score} талбартай бөгөөд нэг илэрсэн нүүрний танилтын үр дүнг илэрхийлнэ.
	\item Сешний түвшин: \texttt{process} командыг нэг удаа ажиллуулахад үүссэн CSV файлыг нэг сешн гэж үзнэ.
	\item Хичээлийн цагийн түвшин: \texttt{periodic} горимд тухайн хичээлийн цагийн \texttt{result\_csv} файлуудыг нэгтгэн \texttt{finalized\_attendance.csv} үүсгэнэ.
	\item Оюутны түвшин: Эцсийн үнэлгээнд оюутан тус бүрийг гараар бүртгэсэн ирцтэй тулган TP, FP, FN (шаардлагатай бол TN)-ийг тооцно.
\end{itemize}

Ингэснээр системийн гүйцэтгэлийг зөвхөн нийт хувь хэмжээгээр бус, хичээлийн цаг тус бүрийн түвшинд задлан үнэлэх боломж бүрдэнэ.

\subsection{Туршилтын протокол ба үнэлгээний алхам}
Үр дүнг бодит хэрэглээнд ойр нөхцөлд тайлагнахын тулд нэгэн жигд протокол ашиглана. Үүнд:

\begin{enumerate}
	\item Хичээлийн цаг бүрийн эцэст системээс үүссэн ирцийн CSV-г хадгална.
	\item Тухайн цагийн гараар бүртгэсэн ирцийн жагсаалтыг лавлах үнэн утга гэж авна.
	\item Системийн ба гараар бүртгэсэн ирцийг оюутан тус бүрээр тулгаж TP, FP, FN, TN ангилалд хуваарилна.
	\item Танигдаагүй (unknown) гэж тэмдэглэгдсэн тохиолдлуудыг тусад нь тоолж, буруу танилттай (FP) хамтад нь шинжилнэ.
	\item Хичээлийн цаг тус бүрийн үзүүлэлтийг бодсоны дараа өдөр/долоо хоногийн нэгтгэл гаргана.
\end{enumerate}

Энэхүү протокол нь нэг өдөрт цөөн сешнтэй байсан ч, олон өдрийн өгөгдөл нэмэгдэхэд ижил зарчмаар үргэлжлүүлэн үнэлгээ хийх боломжтойгоороо давуу талтай.

\subsection{Үнэлгээний хэмжүүрүүдийн хэрэглээ}
Энэ бүлэгт хэмжүүрүүдийн онолын тодорхойлолтыг давтахгүй бөгөөд тухайн туршилтад хэрхэн хэрэглэснийг тайлагнана. TP, FP, TN, FN болон Accuracy/Precision/Recall/F1-ийн тодорхойлолт, томьёог \textit{Судалгааны сэдвийн онол, өнөөгийн түвшин} бүлгийн \textit{Үнэлгээний хэмжүүрүүдийн онолын үндэс} хэсэгт өгсөн.

Үр дүнгийн тайлагналд дараах зарчмыг баримтална:
\begin{itemize}
	\item Арга тус бүрийн ангиллын чанарыг TP, FP, TN, FN бүтэц болон Accuracy-гаар харьцуулна.
	\item Танигдаагүй (unknown) төлвийн тоог тусад нь тайлагнаж, FP/FN-тэй хамтатган trade-off шинжилгээ хийнэ.
	\item Хичээлийн цаг тус бүрийн үр дүнг нэгтгэн өдөр/долоо хоногийн түвшний чиг хандлагыг гаргана.
\end{itemize}

Мөн системийн практик хэрэглээний чадамжийг хугацааны үзүүлэлтүүдээр үнэлнэ. Үүнд сургалтын хугацаа, нэг зураг боловсруулах дундаж хугацаа, нэг нүүр таних дундаж хугацаа, хичээлийн мөчлөгийн нийт хугацаа зэрэг багтана.

\subsection{Үр дүн}

\begin{table}[H]
	\centering
	\caption{Системийн ажиллагаа ба боловсруулах хугацааны бодит үзүүлэлт}
	\label{tab:attendance-operational-template}
	\begin{tabular}{|l|c|}
		\hline
		\textbf{Үзүүлэлт} & \textbf{Утга} \\
		\hline
		Галерейг дахин байгуулах хугацаа (сек) & 460.10 \\
		\hline
		Туршилтын багц дээрх MTCNN танилтын хугацаа (сек) & 67.50 \\
		\hline
		Туршилтын багц дээрх RetinaFace танилтын хугацаа (сек) & 494.44 \\
		\hline
		MTCNN: нэг зураг боловсруулах дундаж хугацаа (секунд) & 1.21 \\
		\hline
		RetinaFace: нэг зураг боловсруулах дундаж хугацаа (секунд) & 8.83 \\
		\hline
		MTCNN: нэг нүүрний мөр боловсруулах дундаж хугацаа (мс) & 259 \\
		\hline
		RetinaFace: нэг нүүрний мөр боловсруулах дундаж хугацаа (мс) & 1653.66 \\
		\hline
	\end{tabular}
\end{table}

\begin{table}[H]
	\centering
	\caption{Leave-one-out validation-ийн бодит үр дүн}
	\label{tab:method-confusion-metrics-template}
	\begin{tabular}{|l|c|c|c|c|c|c|}
		\hline
		\textbf{Арга} & \textbf{N} & \textbf{Closed-set зөв} & \textbf{Closed acc (\%)} & \textbf{Gated зөв} & \textbf{Gated acc (\%)} & \textbf{95\% CI} \\
		\hline
		Cosine & 546 & 539 & 98.72 & 519 & 95.05 & [92.90, 96.58] \\
		\hline
		k-NN & 546 & 538 & 98.53 & 537 & 98.35 & [96.90, 99.13] \\
		\hline
		Hybrid & 546 & 539 & 98.72 & 537 & 98.35 & [96.90, 99.13] \\
		\hline
		SVM & 546 & 538 & 98.53 & 215 & 39.38 & [35.37, 43.54] \\
		\hline
	\end{tabular}
\end{table}

Эндээс харахад hybrid болон k-NN арга нь enrollment өгөгдөл дээр хамгийн тогтвортой байв. Gated accuracy-гаар hybrid/k-NN хоёулаа 98.35\% хүрсэн бөгөөд Cosine нь threshold gating-д илүү мэдрэмтгий байв.

SVM-ийн 39.38\% gated accuracy нь загвар өөрөө бүрэн буруу ажилласнаас илүү open-set gating-ийн тохиргоотой холбоотой. Closed-set зөв ангилалт 98.53\% байсан ч тухайн хувилбарт хэрэглэсэн confidence gate болон unknown шүүлт хатуу байсан тул олон sample эцсийн шийдвэрийн шатанд unknown руу шилжсэн. Өөрөөр хэлбэл, энэ дүн нь SVM-ийн 
``ангилах чадвар''-аас гадна ``шийдвэрийн босго + unknown бодлого''-ын нөлөөг хамтад нь тусгасан үзүүлэлт юм.

Энд тайлагнасан 98.35\% нь enrollment leave-one-out validation-ийн үзүүлэлт бөгөөд бодит анги дахь эцсийн ирцийн үнэн зөвийн шууд хэмжүүр биш гэдгийг тусад нь тэмдэглэх шаардлагатай. Бодит ангийн attendance accuracy, FAR, FNR-ийг бүрэн тооцоход paired ground truth архив зайлшгүй хэрэгтэй.

\subsection{Хичээлийн цаг тус бүрийн ирцийн төлөвийн харьцуулалт}

Энэ хэсэгт хичээлийн цаг тус бүрийн \textit{present/no/unknown} төлөвийн тархалтыг харьцуулсан бөгөөд TP/FP/FN/Precision/Recall зэрэг бүрэн ангиллын хэмжүүрийг тооцоогүй. Учир нь тухайн сешн бүртэй хослох гараар баталгаажуулсан ground truth файл бүрэн хадгалагдаагүй.

\begin{table}[H]
	\centering
	\caption{Хичээлийн цаг ба гэрэлтүүлгийн нөхцөл}
	\label{tab:lesson-metrics-lighting}
	\begin{tabular}{|l|l|l|}
		\hline
		\textbf{Хичээл} & \textbf{Цаг} & \textbf{Гэрэлтүүлэг} \\
		\hline
		1-р паар & 08:00--09:30 & Сүүдэртэй \\
		\hline
		2-р паар & 09:35--11:05 & Сүүдэртэй \\
		\hline
		3-р паар & 11:10--12:40 & Бага зэрэг нарны гэрэлтэй \\
		\hline
		4-р паар & 13:20--14:50 & Үдээс хойш, нарны тусгал нэмэгдэнэ \\
		\hline
		5-р паар & 14:55--16:25 & Үдээс хойш, нарны тусгал ихэснэ \\
		\hline
	\end{tabular}
\end{table}

\begin{table}[H]
	\centering
	\caption{Сешн тус бүрийн баталгаажсан ирцийн нэгтгэл}
	\label{tab:lesson-metrics-by-period}
	\begin{tabular}{|l|l|c|c|c|c|}
		\hline
		\textbf{Өдөр} & \textbf{Хичээл} & \textbf{Мөр} & \textbf{Present} & \textbf{No} & \textbf{Unknown} \\
		\hline
		2026-03-30 & 2nd lesson & 4 & 0 & 4 & 0 \\
		\hline
		2026-03-30 & 4th lesson & 8 & 0 & 8 & 0 \\
		\hline
		2026-03-30 & outside\_period & 16 & 0 & 16 & 0 \\
		\hline
		2026-04-09 & 3rd lesson & 19 & 19 & 0 & 0 \\
		\hline
		2026-04-09 & outside\_period & 14 & 14 & 0 & 0 \\
		\hline
		\textbf{Нийт} & Бүгд & 61 & 33 & 28 & 0 \\
		\hline
		\multicolumn{6}{|l|}{\footnotesize Тайлбар: Мөр=0 байсан 4 сешнийг хүснэгтээс хасаж, дүгнэлтэд нөлөөлөхгүйгээр тусгав.} \\
		\hline
	\end{tabular}
\end{table}

\begin{table}[H]
	\centering
	\caption{Гэрэлтүүлгийн бүсээр нэгтгэсэн ирцийн төлөвийн харьцуулалт}
	\label{tab:lighting-zone-metrics}
	\begin{tabular}{|l|l|c|c|c|c|}
		\hline
		\textbf{Бүс} & \textbf{Хамрах хичээл} & \textbf{Мөр} & \textbf{Present} & \textbf{No} & \textbf{Present rate (\%)} \\
		\hline
		Өглөөний тогтвортой гэрэл & 2-р паар & 4 & 0 & 4 & 0.0 \\
		\hline
		Өдрийн шилжилтийн гэрэл & 3-р паар & 19 & 19 & 0 & 100.0 \\
		\hline
		Үдээс хойших нарны тусгалтай & 4-р паар & 8 & 0 & 8 & 0.0 \\
		\hline
		Хичээлээс гадуурх зураг & outside\_period & 30 & 14 & 16 & 46.7 \\
		\hline
	\end{tabular}
\end{table}

Хүснэгт~\ref{tab:lighting-zone-metrics}-ийн хамгийн онцгой үр дүн нь 3-р паарт present rate 100.0\% байхад 2-р болон 4-р паарт 0.0\% гарсан огцом ялгаа юм. Энэ зөрүүг зөвхөн нэг хүчин зүйлээр тайлбарлах боломжгүй бөгөөд дор хаяж гурван шалтгаан давхцсан гэж үзэв: (1) 3-р паарын сешнд ашиглагдсан зурагнуудад нүүрний чанар, байрлал харьцангуй тогтвортой байсан, (2) 2-р болон 4-р паарын мөрийн тоо цөөн тул босго давсан тогтвортой hit цуглуулахад хангалтгүй байсан, (3) тухайн тохиргооны conservative unknown шүүлт 2-р ба 4-р паарт илүү хүчтэй үйлчилсэн. Иймээс энэ anomaly нь ``систем төгс ажилласан/бүрэн ажиллаагүй'' гэсэн хоёр туйлт дүгнэлтээс илүү, орчны нөхцөл + sampling density + threshold бодлогын хам нөлөө гэж үзэх нь зөв.

% \subsection{Pilot үр дүнг тайлбарлах зарчим}

% \begin{itemize}
% 	\item Арга хоорондын харьцуулалт: Cosine, k-NN, Hybrid (шаардлагатай бол SVM)-ийн TP, FP, FN бүтэц хэрхэн өөр байгааг тайлбарлана.
% 	\item Танигдаагүй төлөвийн нөлөө: Танигдаагүй тохиолдол өсөхөд FP буурах боловч FN өсөх боломжтой солбицлыг дүгнэнэ.
% 	\item Хичээлийн цагийн ялгаа: 1--5-р паарын гэрэл, суудлын өөрчлөлт, хөдөлгөөний нөлөөгөөр үзүүлэлтүүд хэрхэн хэлбэлзэж буйг тайлбарлана.
% \end{itemize}

% Ингэж тайлбарласнаар хүснэгтийн тоон үр дүнг системийн бодит хэрэглээний нөхцөлтэй уялдуулан ойлгуулах боломжтой.

\subsection{Эцсийн ирцийн нэгтгэлийн тайлбар}
Архивлагдсан өгөгдөлд гараар баталгаажуулсан ground truth файл бүрэн хадгалагдаагүй тул ``систем vs гараар бүртгэл'' гэсэн бүрэн comparison-ийг энэ хувилбарт шууд тооцоогүй. Иймээс сешнүүдийн эцсийн ирцийн нэгтгэлийг давхар хүснэгтээр давтахгүйгээр Хүснэгт~\ref{tab:lesson-metrics-by-period}-ийн үр дүнгээр тайлбарласан.

Өөрөөр хэлбэл, энэ бүлгийн attendance нэгтгэл нь \textit{системийн гарцын тайлан} бөгөөд \textit{гараар бүртгэсэн лавлах үнэн утгатай харьцуулсан эцсийн үнэлгээ} биш гэдгийг тодорхой ялгаж ойлгох шаардлагатай.

\subsection{Туршилтын зурагт MTCNN ба RetinaFace-ийн харьцуулалт}
Ижил танилтын боловсруулалтын урсгалыг 56 зурагтай туршилтын багц дээр ашиглаж MTCNN болон RetinaFace илрүүлэгчийг тус тусад нь ажиллуулж хугацаа болон гарцын үзүүлэлтүүдийг харьцууллаа.

\begin{table}[H]
	\centering
	\caption{56 зурагтай туршилтын багц дээрх илрүүлэгчдийн танилтын гүйцэтгэлийн харьцуулалт}
	\label{tab:flat-detector-runtime-comparison}
	\begin{tabular}{|l|c|c|}
		\hline
		\textbf{Үзүүлэлт} & \textbf{MTCNN} & \textbf{RetinaFace} \\
		\hline
		Боловсруулсан зураг & 56 & 56 \\
		\hline
		Нийт runtime (сек) & 67.50 & 494.44 \\
		\hline
		Нэг зураг боловсруулах дундаж хугацаа (сек) & 1.21 & 8.83 \\
		\hline
		Үүссэн face row & 261 & 299 \\
		\hline
		Нэг face row боловсруулах дундаж хугацаа (мс) & 259 & 1653.66 \\
		\hline
		Танигдсан label (known) & 20 & 20 \\
		\hline
		\texttt{unknown} label & 240 & 279 \\
		\hline
		\texttt{no\_face} & 1 & 0 \\
		\hline
		Гаралтын session файл & \texttt{review\_fast\_mtcnn\_2026-04-19\_19-06-03.csv} & \texttt{hybrid\_overhead\_precision\_2026-04-19\_19-44-48.csv} \\
		\hline
	\end{tabular}
\end{table}

Дээрх хүснэгтээс periodic үнэлгээний өнцгөөс хоёр detector-ыг харьцуулахад дараах дүгнэлт гарч байна. Нэгдүгээрт, танигдсан label-ийн тоо ижил (20 vs 20) боловч \texttt{unknown} row-ийн хэмжээ MTCNN дээр 91.95\%, RetinaFace дээр 93.31\% гарсан. Энэ нь туршилтын тухайн профайл дээр detector-оос илүү шийдвэрийн босго болон review-only тохиргоо гаралтын unknown хэмжээг давамгайлан удирдаж байгааг харуулна.

Хоёрдугаарт, 240/261 unknown (91.95\%) үзүүлэлт нь бодит attendance шийдвэр гаргахад шууд ашиглахад өндөр эрсдэлтэй түвшин бөгөөд үүнийг ``систем бүрэн бэлэн'' гэж тайлбарлах боломжгүй. Энэ хувилбарт ашигласан review profile нь зориуд conservative тул FP-ийг дарах оролдлого хийж unknown-ийг өсгөсөн гэж үзсэн. Иймд энэ үр дүнгийн зөв хэрэглээ нь: (1) detector/pipe-ийн харьцуулалт хийх, (2) дараагийн threshold tuning-ийн чиглэл тогтоох, харин (3) эцсийн attendance accuracy гэж шууд дүгнэх биш юм.

Гуравдугаарт, periodic горимд кадр бүрийг миллисекундийн түвшинд боловсруулах шаардлага харьцангуй бага, харин олон мөчлөгийн турш тогтвортой шийдвэр гаргах чадвар чухал. Энэ тохиргоонд MTCNN нь ижил known count-ийг хадгалж, RetinaFace-аас бага unknown шуугиантай гарсан тул pipeline-ийн суурь detector болгон сонгосон.

\subsection{Босгын туршилтын хичээлийн төрлөөрх нэгтгэл}
5 тусдаа босгын туршилтын (0.66, 0.70, 0.74, 0.78, 0.82) сешний CSV-ээс хичээлийн төрлөөр нь дүнг нэгтгэв. Төрөл бүрийн ``хамгийн тохиромжтой босго''-ыг танигдсан давтагдахгүй шошго олон, танигдсан хувь өндөр, \texttt{unknown} хувь бага байх шалгуураар сонгов.

\begin{table}[H]
	\centering
	\caption{Хичээлийн төрөл бүрийн босгын туршилтаас сонгосон хамгийн тохиромжтой утга (MTCNN)}
	\label{tab:threshold-best-by-category-mtcnn}
	\begin{tabular}{|l|c|c|c|c|c|}
		\hline
		\textbf{Төрөл} & \textbf{Сайн босго} & \textbf{Мөр} & \textbf{Танигдсан} & \textbf{Unknown} & \textbf{Танигдсан хувь (\%)} \\
		\hline
		2-р паар & 0.66 & 48 & 4 & 44 & 8.33 \\
		\hline
		3-р паар & 0.66 & 55 & 5 & 50 & 9.09 \\
		\hline
		4-р паар & 0.66 & 2 & 0 & 2 & 0.00 \\
		\hline
		Хичээлээс гадуурх үе & 0.66 & 77 & 3 & 74 & 3.90 \\
		\hline
		Бусад зураг & 0.66 & 35 & 8 & 26 & 22.86 \\
		\hline
	\end{tabular}
\end{table}

Энэ sweep-ийн үр дүнгээр category бүр дээр 0.66 хамгийн ашигтай утга болж гарсан бөгөөд threshold өсөх тусам known rate буурч, unknown rate нэмэгдсэн. Иймд periodic pipeline-ийн дараагийн туршилтад 0.66-ийг baseline болгон авч, required\_hits болон min\_gap параметрүүдтэй хамт цогцоор нь дахин тааруулах нь зүйтэй.

\subsection{Алдааны шинжилгээ}
Системийн алдааг зөвхөн хувь хэмжээгээр дүгнэхээс гадна шалтгаан-үр дагаврын түвшинд задлан шинжлэх шаардлагатай. Практик орчинд FP болон FN алдаанд дараах хүчин зүйлс түгээмэл нөлөөлдөг.

\begin{itemize}
	\item Гэрэлтүүлгийн жигд бус байдал: Арын гэрэл хүчтэй үед нүүрний деталь алдагдаж, эмбеддингийн чанар буурна.
	\item Өнцөг ба байрлалын өөрчлөлт: Камер руу шууд хараагүй нүүрнүүдэд тэгшитгэлийн алдаа ихэснэ.
	\item Хөдөлгөөний бүдгэрэл: Хичээл эхлэх/дуусах үеийн хөдөлгөөн ихтэй мөчид илрүүлэлтийн итгэлцүүр буурах хандлагатай.
	\item Хэсэгчилсэн халхлалт: Маск, үс, гар, дэвтэр зэрэг халхлалт нь танилтын тогтвортой байдалд сөргөөр нөлөөлнө.
\end{itemize}

\begin{figure}[H]
	\centering
	\begin{minipage}[t]{0.32\textwidth}
		\centering
		\includegraphics[width=\linewidth]{../Pictures/TiltedAngleWIthSunlight.png}
		\small (a) Хазайлт + хүчтэй гэрэл: \texttt{танигдаагүй}
	\end{minipage}
	\hfill
	\begin{minipage}[t]{0.32\textwidth}
		\centering
		\includegraphics[width=\linewidth]{../Pictures/TiltedAngle.png}
		\small (b) Ижил төстэй хазайлт: зөв танилт
	\end{minipage}
	\hfill
	\begin{minipage}[t]{0.32\textwidth}
		\centering
		\includegraphics[width=\linewidth]{../Pictures/FalsePositive.png}
		\small (c) FP: өөр хүнийг \texttt{s21c087b\_P.Batsuuri} гэж андуурсан
	\end{minipage}
	\caption{Алдааны кейсийн аннотацтай жишээнүүд}
	\label{fig:error-cases-annotated}
\end{figure}

Туршилтын аннотацтай жишээнүүдээс дараах хоёр төлөөлөх алдаа ажиглагдсан. Нэгдүгээрт, ижил төстэй хазайлттай өнцөг дээр танилт хэвийн байсан боловч нарны шууд тусгал нэмэгдэхэд ижил төстэй байрлалтай нүүрний итгэлцүүр мэдэгдэхүйц буурч, \texttt{танигдаагүй} (\texttt{unknown}) ангилал руу шилжсэн. Энэ нь өнцгийн нөлөөнөөс илүүтэй гэрэлтүүлгийн огцом өөрчлөлт эмбеддингийн тогтвортой байдалд хүчтэй нөлөөлж байгааг харуулж байна. Хоёрдугаарт, FP тохиолдолд бодитоор өөр оюутныг \texttt{s21c087b\_P.Batsuuri} гэж андуурсан нь ажиглагдсан. Энэ төрлийн алдаа нь хоёр оюутны эмбеддингийн зай ойртох, мөн босго параметр харьцангуй сул тохируулагдсан үед өсөх хандлагатай байдаг.

Нэмж хэлэхэд, энэ test image set дээрх MTCNN recognition pass-д \texttt{unknown} label маш өндөр (240/261) гарсан нь review profile нь зориуд conservative тохируулгатай байгааг харуулж байна. Энэ нь manual review-д хэрэглэгдэх үед ашигтай боловч эцсийн attendance баталгаажуулалтад шууд ашиглахын өмнө thresholds-ийг дахин тааруулах шаардлагатай.

Иймээс дараагийн шатанд гэрэлтүүлгийн нөхцөлд мэдрэмтгий байдлыг бууруулах зорилгоор сургалтын өгөгдөлд гэрэл болон өнцгийн төрөлжилт нэмэх, мөн FP бууруулахын тулд шийдвэрийн босго болон зөрүүний босгыг хичээлийн цаг тус бүрээр дахин тохируулах шаардлагатай гэж дүгнэв.

Эдгээр тохиолдлыг аннотацтай зураг болон сешн бүрийн CSV-тэй тулган тайлбарлах нь цаашдын сайжруулалтын шийдвэр гаргалтад шууд ашиг тустай.

\subsection{Хязгаарлалт ба цаашдын ажил}
Энэ туршилтын гол хязгаарлалт нь ангид суурилсан ground truth-ийг тусдаа архивласан CSV хэлбэрээр бүрэн хадгалаагүй, мөн өгөгдөл 22 сурагчтай нэг ангийн хүрээнд төвлөрсөн явдал юм. Иймээс leave-one-out validation өндөр гарсан ч бодит ангид тооцсон FAR/FRR, confidence interval-ийг илүү том, олон өдрийн ground truth дээр дахин баталгаажуулах шаардлагатай.

Мөн гэрэлтүүлгийн хэлбэлзэл нь маш хүчтэй нөлөөлсөн. 3-р паар болон outside\_period-ийн зарим сешнүүдэд систем өндөр чанартай ажилласан бол 2-р болон 4-р паар дээр conservative тохиргоо илүү их no-present шийдвэр гаргасан. Энэ нь параметрүүдийг зөвхөн нэг удаагийн default утгаар бус, хичээлийн цаг тус бүрээр тусад нь тааруулах хэрэгтэйг харуулна.

Цаашид manual attendance ground truth-ийг бүрэн архивлаж, paired statistical test болон bootstrap confidence interval ашиглан methods хоорондын ялгааг илүү нарийвчлалтай баталгаажуулах нь зүйтэй. Мөн UI прототипийг бодит хэрэглэгчийн туршилтаар үнэлж, privacy/consent flow-ийг системийн салшгүй хэсэг болгох шаардлагатай.
\FloatBarrier
